\section*{Resumen ejecutivo}
\addcontentsline{toc}{section}{Resumen ejecutivo}

Pulso reúne membresías, pagos, asistencia, entrenamiento y progreso en una plataforma web. El prototipo documentado incluye entrenador y miembro. Se consultó a quince participantes: cinco responsables y diez miembros.

Se propone una suscripción mensual (SaaS): Starter, hasta 80 miembros e identidad visual básica, por CRC 15\,000; Pro Gym, hasta 200 y configuración ampliada, por CRC 25\,000; Premium, hasta 500 y mayor configuración, por CRC 45\,000; Enterprise, más de 500, desde CRC 70\,000. Los dos últimos podrían requerir diseño propio, atención dedicada y base de datos independiente; necesitan validación y cotización.

La proyección con Starter y Pro Gym termina con 14 gimnasios. Parte de cinco contrataciones y cubre la operación con una más: un Starter y cinco Pro Gym generan CRC 140\,000 frente a CRC 122\,891 de costos, dejando CRC 17\,109. Conseguir esos clientes y sostener los precios requiere pruebas comerciales.

\Needspace{9\baselineskip}
\section{Justificación}

\subsection{Problema y necesidad identificada}

Cada gimnasio necesita saber quién tiene una membresía vigente, qué pagos
requieren seguimiento y qué entrenamiento corresponde a cada miembro. Cuando
esos datos están en registros separados, darles seguimiento y mantenerlos al
día se vuelve más difícil. Pulso reúne esta información en una sola plataforma,
con acceso según el rol: la administración consulta membresías y pagos; el
entrenador organiza los planes de entrenamiento; y cada miembro revisa su rutina
y progreso. Así, Pulso ofrece una forma más clara y coordinada de gestionar las
tareas diarias del gimnasio.

Los responsables consultados utilizan distintas herramientas: dos mencionaron
cuadernos o formularios, dos Excel o Google Sheets, dos software especializado,
dos sistemas internos y uno WhatsApp; cada persona podía seleccionar varias
opciones. Cuatro de cinco responsables señalaron dificultades con los cobros y
pagos pendientes. Estas respuestas ayudan a definir el problema que atiende
Pulso, aunque solo describen los casos consultados en Pérez Zeledón.

\subsection{Oportunidad detectada}

Una plataforma digital puede facilitar la consulta de información y las tareas
diarias del gimnasio. Para que resulte útil, también debe ser fácil de aprender.
Por eso Pulso propone pantallas adaptables al celular y la computadora, ayuda
durante la configuración y funciones organizadas por tarea.

La oferta se dirige primero a dueños y administradores, quienes deciden la
contratación. También debe ser útil para entrenadores y miembros, que la usarán
en sus actividades diarias. En la encuesta, nueve de los diez miembros dieron
la calificación máxima de 5 y uno calificó con 4 la idea de reunir pagos,
asistencia, rutina y progreso.

\subsection{Contexto social, económico, tecnológico y ambiental}

En el ámbito social, una gestión ordenada puede ofrecer mayor claridad sobre
rutinas, asistencia y estado de la membresía. Pulso no garantiza resultados
físicos ni reemplaza el acompañamiento profesional; organiza información para
que las personas autorizadas puedan dar seguimiento.

En el ámbito económico, Pulso ayuda a mantener visibles los pagos pendientes,
las membresías próximas a vencer y la actividad de los miembros. Esto permite que
la administración concentre su tiempo en atender el gimnasio en vez de revisar
registros separados. Durante el piloto se medirá el tiempo empleado antes y
después de utilizar la plataforma para comprobar si facilita esas tareas.

En el ámbito tecnológico, el proyecto utiliza una aplicación web, una base de
datos y una API que comunica las pantallas con el servidor. Los permisos
dependen del rol de cada usuario. Los servicios de nube permiten demostrar el
sistema sin comprar servidores. Para atender clientes será necesario ampliar
la capacidad y preparar respaldos y soporte.

En el ámbito ambiental, los registros digitales podrían reducir el uso de
formularios impresos. Todavía falta medir cuánto papel se sustituiría y qué
consumo generan los servicios de nube.

\subsection{Pertinencia del modelo de suscripción}

El servicio requiere alojamiento, mantenimiento, correcciones, copias de
seguridad y soporte continuo. Por ello se eligió una suscripción mensual en lugar
de una venta única. El piloto de CRC 15\,000 responde al rango indicado por la
única persona que seleccionó un monto y reduce el riesgo de probar una solución
nueva. Después del piloto, los planes propuestos buscan financiar la operación y
deberán validarse con posibles clientes. Cuatro
responsables solicitaron ver una demostración antes de elegir una tarifa; por
esta razón, la estrategia comercial comienza mostrando el producto y su
utilidad, no solamente presentando un precio.

\section{Objetivos del plan de negocios}

\subsection{Objetivo general}

Desarrollar Pulso como una plataforma web de gestión administrativa y deportiva
para gimnasios, como solución para la digitalización, organización y control de
sus procesos.

\subsection{Objetivos específicos}

\begin{enumerate}
  \item Analizar las necesidades administrativas y deportivas de las personas vinculadas a gimnasios en una muestra de quince participantes, como insumo para la propuesta de valor de Pulso.

  \item Desarrollar las funciones de Pulso para gestionar membresías, pagos, asistencia, planes de entrenamiento y progreso.

  \item Definir un modelo de suscripción para la prestación del servicio en línea de Pulso.

  \item Evaluar la viabilidad inicial de Pulso mediante la estimación de costos, ingresos y punto de equilibrio.
\end{enumerate}

El trabajo parte de las necesidades consultadas, continúa con el desarrollo de
funciones y plantea una suscripción para cubrir los costos del servicio. La
proyección financiera permite estudiar cuánto tendría que crecer Pulso para
mantenerse. La encuesta orienta estas decisiones dentro de los límites de sus
quince participantes.
